\documentclass[10pt]{article}
\usepackage[margin=1in]{geometry}
\usepackage[T1]{fontenc}
\usepackage{lmodern,microtype,amsmath,amssymb,amsthm}
\theoremstyle{definition}
\newtheorem{assumption}{Assumption}
\begin{document}
\section*{Smooth Nonconvex Exact-Gradient Stationarity}
We consider a differentiable, possibly nonconvex objective $f:E\to\mathbb{R}$ on a real Hilbert space $E$, with $L$-Lipschitz gradient $\nabla f$ for $L>0$. A global lower bound $\ell\in\mathbb{R}$ satisfies $f(x^0)-\ell\leq\Delta$ at the initial point $x^0\in E$, where $\Delta\geq0$. Each oracle call at $x\in E$ returns $\nabla f(x)$. Given $\varepsilon>0$, success requires a measurable output $\bar{x}\in E$ with integrable squared gradient norm and $\mathbb{E}_{\mathbb{P}}\|\nabla f(\bar{x})\|^2\leq\varepsilon^2$, where $\mathbb{P}$ is the joint law of algorithm and oracle randomness. The quantities $(E,L,\Delta,\varepsilon,x^0)$ are public.

The underlying optimization problem is
\begin{equation}
\label{eq:problem}
\inf_{x\in E} f(x),
\end{equation}
with stationarity as the required accuracy criterion. The space $E$ may be infinite dimensional, carries a specified measurable structure, and has the norm induced by its inner product. The admissible objectives and oracle access satisfy the following conditions.

\begin{assumption}[Differentiability and smoothness]
The function $f$ is differentiable everywhere on $E$, and its gradient satisfies
\begin{equation}
\|\nabla f(x)-\nabla f(y)\|\leq L\|x-y\|
\qquad\text{for every }x,y\in E.
\end{equation}
\end{assumption}

\begin{assumption}[Lower bound and initial gap]
There exists $\ell\in\mathbb{R}$ such that
\begin{align}
\ell &\leq f(x) \qquad\text{for every }x\in E,\\
f(x^0)-\ell &\leq\Delta.
\end{align}
The lower bound need not be attained.
\end{assumption}

\begin{assumption}[Exact-gradient access]
At every call index and every query point $x\in E$, the oracle returns $\nabla f(x)\in E$ for every realization of its randomness. Each query costs one oracle call. Any internal algorithm randomness is independent of the oracle randomness.
\end{assumption}

The objective is accessed through this oracle, and queries may be chosen adaptively anywhere in $E$. The algorithm is selected using the public tuple $(E,L,\Delta,\varepsilon,x^0)$; the objective $f$ and the witnessing lower bound $\ell$ are not supplied as public inputs. An externally specified call budget $K\in\mathbb{N}$ may depend on the entire public tuple.

After $K$ calls, the algorithm returns $\bar{x}:\Omega\to E$, where $(\Omega,\mathbb{P})$ is the probability space carrying the joint randomness. The output must be measurable, and the scalar function $\omega\mapsto\|\nabla f(\bar{x}(\omega))\|^2$ must be integrable with respect to $\mathbb{P}$. The required accuracy is
\begin{equation}
\label{eq:accuracy}
\int_{\Omega}\|\nabla f(\bar{x}(\omega))\|^2\,\mathrm{d}\mathbb{P}(\omega)
\leq\varepsilon^2.
\end{equation}
\end{document}
